% Required packages:
% \usepackage{booktabs}
% \usepackage{graphicx}
% \usepackage{longtable}

\begin{table*}[t]
\centering
\caption{Aggregate four-model results across MedQA, MedMCQA, and PubMedQA. Higher answer accuracy, parse coverage, and SPACE coverage are better. Lower false commitment is better.}
\label{tab:aggregate}
\resizebox{\textwidth}{!}{%
\input{paper_assests/tables/table_main_4model_aggregate_clean.tex}
}
\end{table*}

\begin{table*}[t]
\centering
\caption{Accuracy and reliability deltas of the proposed method compared with each model's best direct or chain-of-thought baseline.}
\label{tab:deltas}
\resizebox{\textwidth}{!}{%
% Auto-generated from CSV. Do not edit manually unless needed.
\begin{tabular}{llllllllllllll}
\toprule
model & proposed\_accuracy & baseline\_name & baseline\_accuracy & absolute\_accuracy\_gain & relative\_accuracy\_gain\_percent & proposed\_parse\_coverage & baseline\_parse\_coverage & proposed\_space\_coverage & baseline\_space\_coverage & proposed\_false\_commitment & baseline\_false\_commitment & false\_commitment\_reduction & model\_clean \\
\midrule
Qwen2.5-7B-Instruct & 0.6367 & Chain-of-thought & 0.6003 & 0.0364 & 6.06 & 1 & 0.944 & 0.9997 & 0.002 & 0.1026 & 0.8478 & 0.7452 & Qwen2.5-7B \\
Phi-4-mini-instruct & 0.5767 & Chain-of-thought & 0.4817 & 0.095 & 19.72 & 0.9913 & 0.7923 & 0.992 & 0 & 0.4054 & 0.9838 & 0.5784 & Phi-4-mini \\
DeepSeek-R1-Distill-Qwen-32B & 0.3877 & Direct & 0.169 & 0.2187 & 129.41 & 0.7213 & 0.3557 & 0.443 & 0.001 & 0.0213 & 0.2986 & 0.2773 & DeepSeek-R1-32B \\
biomistral\_biomistral\_7b & 0.1307 & Chain-of-thought & 0.1247 & 0.006 & 4.81 & 0.333 & 0.3203 & 0.004 & 0 & 0 & 0.8287 & 0.8287 & BioMistral-7B \\
\bottomrule
\end{tabular}
}
\end{table*}

\begin{table*}[t]
\centering
\caption{Per-dataset results for the proposed answer-fusion plus HypothesisMed-v3 SPACE method. Each row uses 1,000 examples.}
\label{tab:perdataset}
\resizebox{\textwidth}{!}{%
% Auto-generated from CSV. Do not edit manually unless needed.
\begin{tabular}{llllllllll}
\toprule
dataset & model & n & answer\_accuracy & ci95 & parse\_coverage & space\_coverage & space\_accuracy & false\_commitment & model\_clean \\
\midrule
medmcqa & Qwen2.5-7B-Instruct & 1000 & 0.561 & [0.530, 0.592] & 1 & 1 & 0.635 & 0.0501 & Qwen2.5-7B \\
medmcqa & Phi-4-mini-instruct & 1000 & 0.489 & [0.458, 0.520] & 0.997 & 0.994 & 0.5191 & 0.3465 & Phi-4-mini \\
medmcqa & DeepSeek-R1-Distill-Qwen-32B & 1000 & 0.34 & [0.311, 0.369] & 0.819 & 0.26 & 0.9923 & 0.0522 & DeepSeek-R1-32B \\
medmcqa & biomistral\_biomistral\_7b & 1000 & 0.164 & [0.141, 0.187] & 0.476 & 0.011 & 0 & 0 & BioMistral-7B \\
medqa & Qwen2.5-7B-Instruct & 1000 & 0.596 & [0.566, 0.626] & 1 & 0.999 & 0.4324 & 0.1361 & Qwen2.5-7B \\
medqa & Phi-4-mini-instruct & 1000 & 0.517 & [0.486, 0.548] & 0.99 & 0.982 & 0.4216 & 0.2156 & Phi-4-mini \\
medqa & DeepSeek-R1-Distill-Qwen-32B & 1000 & 0.255 & [0.228, 0.282] & 0.594 & 0.404 & 0.995 & 0.0118 & DeepSeek-R1-32B \\
medqa & biomistral\_biomistral\_7b & 1000 & 0.17 & [0.147, 0.193] & 0.42 & 0.001 & 1 & 0 & BioMistral-7B \\
pubmedqa & Qwen2.5-7B-Instruct & 1000 & 0.753 & [0.726, 0.780] & 1 & 1 & 0.906 & 0.1215 & Qwen2.5-7B \\
pubmedqa & Phi-4-mini-instruct & 1000 & 0.724 & [0.696, 0.752] & 0.987 & 1 & 0.92 & 0.654 & Phi-4-mini \\
pubmedqa & DeepSeek-R1-Distill-Qwen-32B & 1000 & 0.568 & [0.537, 0.599] & 0.751 & 0.665 & 0.997 & 0 & DeepSeek-R1-32B \\
pubmedqa & biomistral\_biomistral\_7b & 1000 & 0.058 & [0.044, 0.072] & 0.103 & 0 &  & 0 & BioMistral-7B \\
\bottomrule
\end{tabular}
}
\end{table*}

\begin{table}[t]
\centering
\caption{Confidence and accuracy summary by dataset for the proposed fusion outputs.}
\label{tab:confidence_dataset}
\resizebox{\linewidth}{!}{%
% Auto-generated from CSV. Do not edit manually unless needed.
\begin{tabular}{llll}
\toprule
dataset & mean\_confidence & accuracy & n \\
\midrule
medmcqa & 0.1208 & 0.3885 & 4000 \\
medqa & 0.1282 & 0.3845 & 4000 \\
pubmedqa & 0.2201 & 0.5258 & 4000 \\
\bottomrule
\end{tabular}
}
\end{table}

% Optional appendix table. Use only if you want example-level confidence rows.
% For this one, add \usepackage{longtable}.
\begin{table*}[t]
\centering
\caption{Example-level confidence rows from proposed fusion outputs, truncated to the first 30 rows for readability.}
\label{tab:confidence_rows}
\resizebox{\textwidth}{!}{%
% Auto-generated from CSV. Do not edit manually unless needed.
\begin{tabular}{lllll}
\toprule
file & model & dataset & confidence & correct \\
\midrule
results/fusion/biomistral\_biomistral\_7b\_fusion\_majority\_answer\_hypmed\_v3\_space\_medmcqa\_original1000.jsonl & biomistral\_biomistral\_7b & medmcqa & 0 & 0 \\
results/fusion/biomistral\_biomistral\_7b\_fusion\_majority\_answer\_hypmed\_v3\_space\_medmcqa\_original1000.jsonl & biomistral\_biomistral\_7b & medmcqa & 0 & 0 \\
results/fusion/biomistral\_biomistral\_7b\_fusion\_majority\_answer\_hypmed\_v3\_space\_medmcqa\_original1000.jsonl & biomistral\_biomistral\_7b & medmcqa & 0 & 0 \\
results/fusion/biomistral\_biomistral\_7b\_fusion\_majority\_answer\_hypmed\_v3\_space\_medmcqa\_original1000.jsonl & biomistral\_biomistral\_7b & medmcqa & 0 & 0 \\
results/fusion/biomistral\_biomistral\_7b\_fusion\_majority\_answer\_hypmed\_v3\_space\_medmcqa\_original1000.jsonl & biomistral\_biomistral\_7b & medmcqa & 0 & 0 \\
results/fusion/biomistral\_biomistral\_7b\_fusion\_majority\_answer\_hypmed\_v3\_space\_medmcqa\_original1000.jsonl & biomistral\_biomistral\_7b & medmcqa & 0 & 0 \\
results/fusion/biomistral\_biomistral\_7b\_fusion\_majority\_answer\_hypmed\_v3\_space\_medmcqa\_original1000.jsonl & biomistral\_biomistral\_7b & medmcqa & 0 & 1 \\
results/fusion/biomistral\_biomistral\_7b\_fusion\_majority\_answer\_hypmed\_v3\_space\_medmcqa\_original1000.jsonl & biomistral\_biomistral\_7b & medmcqa & 0 & 0 \\
results/fusion/biomistral\_biomistral\_7b\_fusion\_majority\_answer\_hypmed\_v3\_space\_medmcqa\_original1000.jsonl & biomistral\_biomistral\_7b & medmcqa & 0 & 0 \\
results/fusion/biomistral\_biomistral\_7b\_fusion\_majority\_answer\_hypmed\_v3\_space\_medmcqa\_original1000.jsonl & biomistral\_biomistral\_7b & medmcqa & 0 & 0 \\
results/fusion/biomistral\_biomistral\_7b\_fusion\_majority\_answer\_hypmed\_v3\_space\_medmcqa\_original1000.jsonl & biomistral\_biomistral\_7b & medmcqa & 0 & 0 \\
results/fusion/biomistral\_biomistral\_7b\_fusion\_majority\_answer\_hypmed\_v3\_space\_medmcqa\_original1000.jsonl & biomistral\_biomistral\_7b & medmcqa & 0 & 0 \\
results/fusion/biomistral\_biomistral\_7b\_fusion\_majority\_answer\_hypmed\_v3\_space\_medmcqa\_original1000.jsonl & biomistral\_biomistral\_7b & medmcqa & 0 & 1 \\
results/fusion/biomistral\_biomistral\_7b\_fusion\_majority\_answer\_hypmed\_v3\_space\_medmcqa\_original1000.jsonl & biomistral\_biomistral\_7b & medmcqa & 0 & 0 \\
results/fusion/biomistral\_biomistral\_7b\_fusion\_majority\_answer\_hypmed\_v3\_space\_medmcqa\_original1000.jsonl & biomistral\_biomistral\_7b & medmcqa & 0 & 0 \\
results/fusion/biomistral\_biomistral\_7b\_fusion\_majority\_answer\_hypmed\_v3\_space\_medmcqa\_original1000.jsonl & biomistral\_biomistral\_7b & medmcqa & 0 & 0 \\
results/fusion/biomistral\_biomistral\_7b\_fusion\_majority\_answer\_hypmed\_v3\_space\_medmcqa\_original1000.jsonl & biomistral\_biomistral\_7b & medmcqa & 0 & 0 \\
results/fusion/biomistral\_biomistral\_7b\_fusion\_majority\_answer\_hypmed\_v3\_space\_medmcqa\_original1000.jsonl & biomistral\_biomistral\_7b & medmcqa & 0 & 1 \\
results/fusion/biomistral\_biomistral\_7b\_fusion\_majority\_answer\_hypmed\_v3\_space\_medmcqa\_original1000.jsonl & biomistral\_biomistral\_7b & medmcqa & 0 & 0 \\
results/fusion/biomistral\_biomistral\_7b\_fusion\_majority\_answer\_hypmed\_v3\_space\_medmcqa\_original1000.jsonl & biomistral\_biomistral\_7b & medmcqa & 0 & 0 \\
results/fusion/biomistral\_biomistral\_7b\_fusion\_majority\_answer\_hypmed\_v3\_space\_medmcqa\_original1000.jsonl & biomistral\_biomistral\_7b & medmcqa & 0 & 0 \\
results/fusion/biomistral\_biomistral\_7b\_fusion\_majority\_answer\_hypmed\_v3\_space\_medmcqa\_original1000.jsonl & biomistral\_biomistral\_7b & medmcqa & 0 & 0 \\
results/fusion/biomistral\_biomistral\_7b\_fusion\_majority\_answer\_hypmed\_v3\_space\_medmcqa\_original1000.jsonl & biomistral\_biomistral\_7b & medmcqa & 0 & 0 \\
results/fusion/biomistral\_biomistral\_7b\_fusion\_majority\_answer\_hypmed\_v3\_space\_medmcqa\_original1000.jsonl & biomistral\_biomistral\_7b & medmcqa & 0 & 0 \\
results/fusion/biomistral\_biomistral\_7b\_fusion\_majority\_answer\_hypmed\_v3\_space\_medmcqa\_original1000.jsonl & biomistral\_biomistral\_7b & medmcqa & 0 & 0 \\
results/fusion/biomistral\_biomistral\_7b\_fusion\_majority\_answer\_hypmed\_v3\_space\_medmcqa\_original1000.jsonl & biomistral\_biomistral\_7b & medmcqa & 0 & 0 \\
results/fusion/biomistral\_biomistral\_7b\_fusion\_majority\_answer\_hypmed\_v3\_space\_medmcqa\_original1000.jsonl & biomistral\_biomistral\_7b & medmcqa & 0 & 0 \\
results/fusion/biomistral\_biomistral\_7b\_fusion\_majority\_answer\_hypmed\_v3\_space\_medmcqa\_original1000.jsonl & biomistral\_biomistral\_7b & medmcqa & 0 & 1 \\
results/fusion/biomistral\_biomistral\_7b\_fusion\_majority\_answer\_hypmed\_v3\_space\_medmcqa\_original1000.jsonl & biomistral\_biomistral\_7b & medmcqa & 0 & 0 \\
results/fusion/biomistral\_biomistral\_7b\_fusion\_majority\_answer\_hypmed\_v3\_space\_medmcqa\_original1000.jsonl & biomistral\_biomistral\_7b & medmcqa & 0 & 1 \\
\bottomrule
\end{tabular}
}
\end{table*}
